\documentclass[11pt]{article}

\usepackage[T1]{fontenc}
\usepackage[utf8]{inputenc}
\usepackage[margin=1in]{geometry}
\usepackage{amsmath,amssymb}
\usepackage{booktabs}
\usepackage{tabularx}
\usepackage{array}
\usepackage{microtype}
\usepackage{natbib}
\usepackage[hidelinks]{hyperref}
\usepackage{url}

\title{Commodity Tail Dependence and Exchange-Rate Risk in Ghana:\\
Copula Adequacy, Stress Tests, and Calendar Robustness}
\author{Jonathan Nelson\\Independent Research Project}
\date{}

\begin{document}

\maketitle

\begin{abstract}
This study examines commodity tail dependence and exchange-rate risk in Ghana using two co-equal calendar constructions for cocoa, gold, Brent crude oil, WTI crude oil and the Ghanaian cedi over 2015--2026. The analysis combines explicit marginal convergence and adequacy gates, eight static bivariate copula families evaluated with a Rosenblatt-transform Cramér--von Mises statistic and parametric bootstrap, finite-threshold lower- and upper-tail concentration at prespecified thresholds, and moving-block-bootstrap stress comparisons corrected for multiple testing. Across the combined COVID-19 plus 2024 stress definition, COVID-19 alone and 2024 alone, no stress-calm comparison survives either Bonferroni or Benjamini--Hochberg correction in either calendar panel. Copula adequacy is more heterogeneous. Gold--Brent, Gold--WTI and Brent--WTI reject all eight tested static families under both calendar constructions, while all four commodity--Cedi pairs reject none of the eight in either panel. Cocoa--Gold, Cocoa--Brent and Cocoa--WTI are calendar-sensitive, with complete rejection under business-day alignment but only partial rejection under common-observation alignment. The Cedi remains inadequately represented by the tested continuous marginal models under both panels, so Cedi-related copula and extreme-value estimates are interpreted as conditional diagnostics. The results show that calendar alignment, absolute model adequacy and multiplicity control materially affect which dependence conclusions are defensible.
\end{abstract}

\noindent\textbf{Keywords:} tail dependence; copulas; commodity risk; Ghana; exchange-rate risk; model adequacy; calendar alignment; stress testing

\section{Introduction}
\label{sec:introduction}

Ghana's external sector remains strongly exposed to commodity markets. UN Trade and Development classifies an economy as commodity-dependent when commodities account for more than 60\% of merchandise exports, and its 2025 country profile reports that Ghana's three leading commodity exports represented 78.4\% of allocated product exports over 2021--2023. Bank of Ghana data show the continuing importance of cocoa, gold and crude oil: in 2025, exports of cocoa beans and products were about US$4.0 billion, non-monetary gold about US$21.0 billion, and crude oil about US$2.6 billion. These exposures make the relation between commodity prices and the Ghanaian cedi relevant not only for macroeconomic analysis but also for financial risk measurement and decisions made under extreme market conditions \citep{unctad2025,bog2025}.

The Ghana-specific literature already shows that this relation is neither simple nor constant. \citet{archer2022} use monthly cocoa, gold, Brent crude oil and exchange-rate data and report asymmetric relationships across exchange-rate quantiles. \citet{barson2023} use wavelet methods and find that the comovement of Ghana's major commodity returns varies across time horizons, with the exchange rate playing a material role in those relationships. \citet{yeboah2025}, using quantile-on-quantile estimation over 2003--2023, likewise report nonlinear and asymmetric relationships among commodity prices, inflation and the exchange rate. Recent work has widened the setting further. \citet{atipaga2025} use daily data to examine the Cedi's connection with global uncertainty measures, while \citet{woode2025} study time- and frequency-varying connectedness among commodities, currencies and equity markets in commodity-dependent sub-Saharan African economies, including Ghana. Together, these studies make a fixed linear commodity-Cedi relationship difficult to justify.

International evidence points in the same direction. Copula studies of commodity currencies find that dependence can vary with the commodity, country, sample period and market state. \citet{reboredo2012} reports generally weak oil-exchange-rate dependence in the currencies studied, with stronger dependence after the global financial crisis. \citet{ignatieva2017} find positive dependence between commodity-price indices and the currencies of major commodity exporters, together with an increase in time-varying tail dependence after the crisis. More recent studies continue to find state dependence. \citet{go2024} document changing return and volatility spillovers between Malaysian crude palm oil and exchange rates, while \citet{dodd2026} show that the usual positive commodity-currency return relation can disappear when political risk is high. Crisis-focused work also provides reasons to examine the tails directly: \citet{qiao2023}, for example, report greater lower- and upper-tail contagion across commodity futures after the COVID-19 outbreak. Such findings motivate direct testing of stress-period tail behavior rather than assuming that dependence must increase in every crisis or market.

Copula methods are therefore well suited to the problem, but their use is not itself new in Ghana. \citet{mensah2020} estimate static and time-varying copulas for major foreign currencies against the Cedi and document tail-dependent exchange-rate comovement. The relevant gap is narrower. Existing Ghana studies have not centered the question of whether commodity-Cedi tail conclusions remain stable when the underlying daily data are aligned using different defensible market-calendar rules. That issue matters because commodity futures and the Cedi do not necessarily share the same trading days or observation timing. The classical nonsynchronous-trading literature shows that mismatched observation times can distort return-based dependence estimates. \citet{martens2001}, for example, show that close-to-close international stock returns can understate correlation and that synchronization changes estimated correlation dynamics. Related work finds that time adjustment can materially alter market-integration and spillover estimates. Calendar construction should therefore be treated as part of the statistical design rather than as an invisible data-cleaning step.

A second issue is model adequacy. Selecting the copula with the best relative information criterion does not establish that the selected family provides an adequate absolute representation of the data. Copula goodness-of-fit research distinguishes these questions and commonly relies on resampling when parameters are estimated. The present study therefore evaluates eight static bivariate copula families using a Rosenblatt-transform Cramér-von Mises statistic with parametric bootstrap, rather than treating relative fit as evidence of correctness. The marginal stage is also subject to explicit convergence and diagnostic gates because misspecified margins can distort copula and downstream risk estimates. This point is particularly relevant for the Cedi, whose return distribution contains many exact zeros and whose selected continuous marginal specification does not pass the same probability-integral-transform adequacy tests as the commodity margins. Cedi-related copula and extreme-value results are consequently interpreted as conditional diagnostics rather than model-free structural quantities.

A third issue concerns inference in the tails. The study uses empirical lower- and upper-tail concentration at prespecified finite thresholds rather than presenting these quantities as asymptotic tail-dependence coefficients. This distinction matters because finite-sample tail-dependence estimates are sensitive to the estimator and threshold. Stress-period comparisons are evaluated with a moving-block bootstrap, and multiplicity is controlled within each prespecified 60-test family using Bonferroni and Benjamini-Hochberg procedures. This design responds to a practical problem in large tail-comparison exercises: a collection of nominal p-values can look persuasive even when no result survives correction. Recent formal work on differences and equivalence between tail copulas, including \citet{bormann2020}, \citet{can2024}, and \citet{koike2025}, reinforces the need to state precisely which tail quantity is being compared and what inferential claim a test can support.

Against this background, the paper uses two co-equal constructions of the same daily-source data. The first preserves a business-day calendar and allows limited forward filling before returns are computed. The second computes close-to-close returns only between consecutive dates on which all five series have actual observations. Neither construction is treated as the truth or as a correction of the other. Instead, substantive findings are classified as calendar-robust when they survive both constructions, calendar-sensitive when the conclusions differ, and panel-specific when they arise only as a diagnostic of one construction. This makes calendar sensitivity an explicit part of the evidence rather than a post hoc robustness exercise.

The analysis addresses three questions. First, do empirical lower- and upper-tail concentration measures differ between calm periods and the COVID-19 and 2024 stress periods, and are those conclusions robust to alternative calendar-alignment methods? Second, can any of eight tested static bivariate copula families adequately represent the dependence structure among Ghana-relevant commodity and exchange-rate pairs, and which adequacy conclusions are robust across calendar constructions? Third, how sensitive are Cedi-related dependence and extreme-value diagnostics to the Cedi's zero-heavy return distribution, marginal-model inadequacy and calendar construction? By answering these questions jointly, the paper shifts attention from selecting a single preferred dependence model to identifying which conclusions remain defensible after calendar construction, marginal adequacy, absolute copula fit and multiple testing are taken seriously.

\section{Related Literature}
\label{sec:related_literature}

\subsection{Commodity prices, exchange rates and connectedness in Ghana}

The relationship between commodity markets and the Ghanaian cedi has been studied from several perspectives, but the literature has developed mostly through monthly-frequency models of asymmetric response, volatility transmission and time-frequency comovement. \citet{zankawah2020}, using monthly data from 1991 to 2015 and multivariate GARCH specifications, report significant spillover effects from crude oil prices to Ghana's exchange rate. Their analysis is important because it treats exchange-rate volatility as part of a broader commodity-shock system, but it does not examine joint tail dependence or the sensitivity of daily dependence estimates to market-calendar construction.

A later group of studies broadened the commodity set and placed greater emphasis on nonlinear relationships. \citet{archer2022} analyse cocoa, gold, Brent crude oil and the nominal exchange rate with quantile regression over September 2007 to December 2020. Their results vary across exchange-rate quantiles, which is consistent with the view that commodity-currency relationships can depend on market conditions rather than follow a single linear pattern. \citet{boateng2022} similarly document time-frequency variation in the relationships among cocoa, Brent crude oil, gold and Ghanaian macroeconomic indicators using wavelet methods. Although their focus is not the Cedi alone, their findings reinforce the broader point that commodity linkages in Ghana vary by horizon and state.

\citet{barson2023} move closer to the present question by studying commodity-return comovement in Ghana and the role of the exchange rate. Using monthly observations from September 2007 to March 2021, they find stronger coherence at short and medium horizons and conclude that the exchange rate plays an important role in commodity interdependence. \citet{yeboah2025} provide more recent evidence using quantile-on-quantile estimation with monthly data from 2003 to 2023. Their results also point to asymmetric and nonlinear relationships between commodity prices and the exchange rate. Taken together, these studies make a constant, symmetric commodity-Cedi relationship difficult to defend.

Recent work has expanded the Ghanaian literature beyond direct commodity-price regressions. \citet{atipaga2025} use daily Cedi data from June 2016 to October 2023 with wavelet and Rényi transfer entropy methods to examine global uncertainty channels. Their results associate Cedi dynamics with oil-price volatility and Ghana's sovereign spread, with effects varying across horizons. \citet{woode2025} examine time- and frequency-varying connectedness among commodities, currencies and equity markets in commodity-dependent sub-Saharan African economies, including Ghana, and report changing transmission patterns across crises and horizons. These studies confirm that frequency, market state and model choice matter for how cross-market dependence is described.

Copula methods themselves are not new to Ghana. \citet{mensah2020} use daily GHS exchange rates against the U.S. dollar, euro and pound sterling from 2009 to 2019 and estimate both static and time-varying copulas. They report positive dependence and evidence that tail dependence evolves over time. This study is an important novelty guardrail for the present paper. The contribution here is not the introduction of copulas to Ghanaian financial data. Rather, the question is whether commodity-Cedi dependence, copula adequacy and stress-period tail conclusions remain defensible under alternative treatments of nonsynchronous market calendars, explicit marginal diagnostics and multiplicity-controlled inference.

The Ghana studies reviewed here therefore establish three points. First, commodity and exchange-rate relationships are economically relevant and empirically non-constant. Second, nonlinear and time-frequency methods often reveal features that are hidden in mean-based specifications. Third, there is already a domestic precedent for copula-based tail analysis. What remains less developed in this literature is the interaction between daily market-calendar construction, absolute copula adequacy, marginal misspecification and corrected inference across many tail comparisons. That narrower gap motivates the design adopted in this study.

\subsection{International commodity-currency dependence and tail risk}

The international commodity-currency literature also shows substantial heterogeneity. \citet{reboredo2012} uses copulas to model oil-price and exchange-rate comovement and finds generally weak dependence in the currencies studied, although dependence strengthened after the global financial crisis. \citet{ignatieva2017}, examining major commodity exporters at daily frequency, find positive dependence between commodity-price indices and commodity currencies and report a pronounced increase in time-varying tail dependence after the global financial crisis. These results are not contradictory so much as evidence that dependence is sensitive to country, commodity composition, sample period and model specification.

Later studies place greater emphasis on extreme risk rather than average comovement. \citet{bondatti2022} analyse downside commodity tail risk in exchange rates and show that currency exposure is heterogeneous across commodities, including oil and gold. \citet{ning2025} use a dependence-switching copula framework to study oil prices and exchange rates in both oil-exporting and oil-importing countries and report that dependence and downside/upside spillovers vary across regimes. Such results caution against treating a commodity currency as having one stable dependence coefficient that applies equally in calm and stressed markets.

Recent evidence extends that argument. \citet{go2024}, using Malaysian crude palm oil and foreign-exchange markets, document changes in return and volatility spillovers across crisis and post-crisis periods. \citet{dodd2026} show that the usually positive contemporaneous relationship between commodity and currency returns in major commodity-exporting countries can disappear when political risk rises. Their result is especially useful for the present paper because it shows that even a well-established commodity-currency relation can be state-dependent. The implication is not that Ghana must follow the same pattern, but that robustness to market conditions should be tested rather than assumed.

This literature also helps define what the present paper does not claim. It does not estimate causal transmission from commodity prices to the Cedi, and it does not attempt to identify a universal commodity-currency pricing mechanism. Instead, it asks whether observed dependence and tail-concentration conclusions are stable across two defensible calendar constructions and whether any of eight static bivariate copula families provides an adequate representation of those relationships. This narrower objective is appropriate given the heterogeneity documented in the broader literature.

\subsection{Crisis periods and changes in extreme dependence}

A large post-2020 literature reports stronger cross-market dependence during periods of financial distress. In commodity markets, \citet{qiao2023} use a copula-based network approach and document higher lower- and upper-tail contagiousness after the onset of COVID-19. Chen et al. (2023) likewise find that dependence among commodity futures changes substantially during major distress episodes when flexible dynamic copula models are used. Related quantile-connectedness work also reports stronger transmission in the tails during the pandemic.

These findings provide a strong reason to test stress-period changes, but they do not establish that every commodity pair or every currency relationship must show stronger tail dependence in a crisis. Samples differ in market coverage, frequency, stress-window definition, tail estimand and model structure. They also differ in how many hypotheses are evaluated and whether multiplicity is controlled. The present study therefore treats COVID-19 and 2024 as prespecified stress windows and asks whether finite-threshold lower- and upper-tail concentration differs from the calm baseline under the same inferential design in both calendar panels.

This distinction is important because nominal significance can accumulate quickly in a large set of pair, threshold and tail comparisons. \citet{bormann2020} provide a close methodological precedent in showing that structural differences in financial tail dependence can be studied with bootstrap procedures while accounting for multiple comparisons. More recent work develops formal two-sample tests for equality of tail copulas (Can, Einmahl and Laeven, 2024) and tests of tail equivalence (Koike, Kato and Yoshiba, 2025). These procedures address broader inferential objects than the finite-threshold concentration differences used here, but they reinforce a common principle: the statistical claim must match the tail quantity actually tested.

Accordingly, the present stress analysis is intentionally modest in interpretation. A failure to reject a stress-calm difference is not evidence that the two regimes are identical, and an unadjusted p-value below 0.05 is not treated as a robust discovery when the family contains many simultaneous tests. The relevant question is whether a result survives the prespecified multiplicity correction and whether that corrected conclusion agrees across both calendar constructions.

\subsection{Nonsynchronous trading and calendar construction}

The market-calendar problem has a long history in empirical finance. \citet{scholes1977} show that nonsynchronous trading can introduce serious errors-in-variables problems into return-based models. \citet{lo1990} formalize how infrequent or nonsynchronous trading affects variances, autocorrelations and cross-autocorrelations. These early studies establish that observed return dependence can reflect the observation process as well as the underlying economic relationship.

The problem becomes particularly visible when assets trade in different markets or at different times. \citet{martens2001} show that close-to-close international stock returns can understate correlation when trading hours differ and that estimated covariance and correlation dynamics change after synchronization. \citet{schotman2006}, studying Central European emerging markets, similarly find that time adjustment improves market-integration estimates and that aggregation to weekly frequency does not eliminate the time-matching problem. More recently, \citet{dimitriou2025} show that synchronization choices materially affect volatility-spillover inference in G-7 equity markets during crisis periods.

These studies do not prescribe one universally correct calendar treatment for the present data. Their contribution is more fundamental: calendar alignment is part of model construction and can change measured dependence. That is why this study uses two co-equal panels rather than designating one as the true dataset and the other as a robustness check. The business-day panel preserves the documented business-day construction with limited forward filling, while the common-observation panel computes returns only between consecutive dates on which all five series have actual observations. The resulting comparison asks which conclusions survive both plausible representations of the observation process.

The common-observation panel should not be interpreted as perfect intraday synchronization. Its return intervals range across consecutive common observation dates and can span more than one calendar day. The business-day panel, by contrast, preserves nominal business-day frequency but can generate zero returns when prices are carried forward. These are different compromises. Treating the two constructions symmetrically makes the consequences of those compromises visible rather than embedding them silently in a single preprocessing choice.

\subsection{Copula adequacy, marginal specification and tail inference}

Copulas separate marginal behavior from cross-variable dependence, which makes them attractive for financial returns with non-Gaussian margins. That separation does not make the marginal stage irrelevant. \citet{fantazzini2009} shows through simulation that misspecified margins can materially bias estimated copula parameters and risk measures. This is directly relevant to the Cedi in the present study because its return series contains many exact zeros and the selected continuous marginal model fails the same probability-integral-transform adequacy diagnostics passed by the commodity margins. The resulting Cedi-related copula and extreme-value quantities are therefore interpreted conditionally on an inadequate marginal reference model rather than as model-free structural facts.

The treatment of ties deserves similar care. \citet{kojadinovic2017} shows that many standard copula procedures are derived under the assumption of continuous margins without ties and can be biased when ties arise through rounding or limited measurement precision. The present project therefore audited ties explicitly. Panel B contains no exact PIT ties. Panel A contains one three-way Cedi PIT tie created by the numerical clipping of three extreme observations at \(10^{-6}\). Recovering their underlying rank order changes only two pseudo-observations, leaves membership in the prespecified \(q=.025,.05,.10\) lower-tail sets unchanged, and produces only negligible changes in the observed Cedi-pair goodness-of-fit statistics. The clipping issue is therefore treated as a closed numerical audit rather than evidence that all Cedi copula inference is invalid.

Goodness of fit also needs to be distinguished from relative model selection. \citet{genestremillard2008} establish conditions under which parametric bootstrap methods provide valid approximations for goodness-of-fit testing in semiparametric models. \citet{genest2009} review a broad class of copula goodness-of-fit procedures and emphasize bootstrap implementation under composite null hypotheses. \citet{dobric2007} study a goodness-of-fit procedure based on the Rosenblatt transformation and show that bootstrap treatment is important when margins are estimated rather than known. Consistent with this literature, the present paper reports a Rosenblatt-transform Cramér-von Mises statistic with parametric bootstrap. It does not equate the lowest AIC with model adequacy, and it does not describe the implemented statistic as an exact replication of every empirical-copula procedure considered by Genest, Rémillard and Beaudoin.

Tail measurement presents a separate problem. \citet{frahm2005} show that estimators of the asymptotic tail-dependence coefficient can behave differently in finite samples, while \citet{supper2020} show that static estimators can be severely biased when extreme dependence varies over time. These results support two wording choices in the present study. First, the empirical quantities calculated at \(q=.025,.05,.10\) are described as finite-threshold tail concentration measures rather than asymptotic tail-dependence coefficients. Second, full-sample static copula estimates are treated as summaries of the observed sample rather than as evidence that the underlying dependence process is time-invariant.

Finally, the large number of stress comparisons requires explicit multiplicity control. \citet{benjamini1995} introduced false-discovery-rate control as an alternative to family-wise error control, while Bonferroni adjustment provides a stricter family-wise criterion. The present analysis reports both. Each stress definition is treated as its own prespecified 60-test family, comprising ten pairs, three thresholds and two tails. This prevents isolated nominal p-values from being interpreted as robust stress effects when they do not survive the correction applied to the family from which they were selected.

\subsection{Positioning of the present study}

The literature reviewed above points to a specific rather than sweeping gap. Ghanaian studies already document asymmetric commodity-exchange-rate relationships, time-frequency connectedness and exchange-rate tail dependence. International studies already demonstrate that commodity-currency dependence can be nonlinear, dynamic and crisis-sensitive. Statistical research already provides formal tools for copula adequacy, tail comparison and multiple testing.

The present study connects these strands through a robustness question that is rarely made explicit in applied commodity-currency work: which conclusions survive alternative defensible treatments of nonsynchronous market calendars? It addresses that question using two co-equal panels, explicit marginal convergence and adequacy gates, eight static bivariate copula families with absolute goodness-of-fit testing, finite-threshold stress comparisons, and Bonferroni and Benjamini-Hochberg correction within prespecified test families. It also separates calendar-robust findings from calendar-sensitive and panel-specific diagnostics. The contribution is therefore not a claim that a new dependence method has been invented or that copulas have not previously been used in Ghana. It is a design for distinguishing conclusions that are stable to consequential data-construction choices from those that are not.

\subsection{Working reference register for this section}

\begin{itemize}
\item Archer, C., Owusu Junior, P., Adam, A. M., Asafo-Adjei, E., \& Baffoe, S. (2022). *Asymmetric Dependence Between Exchange Rate and Commodity Prices in Ghana*. Annals of Financial Economics, 17(2). DOI: 10.1142/S2010495222500129.
\item Atipaga, U.-F., \& Mensah, M. (2025). *Global uncertainties and the Ghanaian cedi: time-frequency and directional information flow analysis*. Cogent Economics \& Finance, 13(1), 2598191. DOI: 10.1080/23322039.2025.2598191.
\item Barson, Z., Owusu Junior, P., \& Adam, A. M. (2023). *Comovement between commodity returns in Ghana: the role of exchange rates*. Journal of Economic Structures, 12, 17. DOI: 10.1186/s40008-023-00312-z.
\item Benjamini, Y., \& Hochberg, Y. (1995). *Controlling the False Discovery Rate: A Practical and Powerful Approach to Multiple Testing*. Journal of the Royal Statistical Society Series B, 57(1), 289--300. DOI: 10.1111/j.2517-6161.1995.tb02031.x.
\item Boateng, E., Asafo-Adjei, E., Addison, A., Quaicoe, S., Yusuf, M. A., Abeka, M. J., \& Adam, A. M. (2022). *Interconnectedness among commodities, the real sector of Ghana and external shocks*. Resources Policy, 75, 102511. DOI: 10.1016/j.resourpol.2021.102511.
\item Bondatti, M., \& Rillo, G. (2022). *Commodity tail-risk and exchange rates*. Finance Research Letters, 47, 102937. DOI: 10.1016/j.frl.2022.102937.
\item Bormann, C., \& Schienle, M. (2020). *Detecting Structural Differences in Tail Dependence of Financial Time Series*. Journal of Business \& Economic Statistics, 38(2), 380--392. DOI: 10.1080/07350015.2018.1506343.
\item Can, S. U., Einmahl, J. H. J., \& Laeven, R. J. A. (2024). *Two-Sample Testing for Tail Copulas with an Application to Equity Indices*. Journal of Business \& Economic Statistics, 42(1), 147--159. DOI: 10.1080/07350015.2023.2166050.
\item Dimitriou, D., \& Kenourgios, D. (2025). *The implications of non-synchronous trading in G-7 financial markets*. International Journal of Finance \& Economics, 30(1), 689--709. DOI: 10.1002/ijfe.2936.
\item Dobrić, J., \& Schmid, F. (2007). *A goodness of fit test for copulas based on Rosenblatt's transformation*. Computational Statistics \& Data Analysis, 51(9), 4633--4642. DOI: 10.1016/j.csda.2006.08.012.
\item Dodd, O., Fernandez-Perez, A., \& Sosvilla-Rivero, S. (2026). *Political risk and commodity currencies*. Journal of Commodity Markets, 42, 100562. DOI: 10.1016/j.jcomm.2026.100562.
\item Fantazzini, D. (2009). *The effects of misspecified marginals and copulas on computing the value at risk: A Monte Carlo study*. Computational Statistics \& Data Analysis, 53(6), 2168--2188. DOI: 10.1016/j.csda.2008.02.002.
\item Frahm, G., Junker, M., \& Schmidt, R. (2005). *Estimating the tail-dependence coefficient: Properties and pitfalls*. Insurance: Mathematics and Economics, 37(1), 80--100. DOI: 10.1016/j.insmatheco.2005.05.008.
\item Genest, C., \& Rémillard, B. (2008). *Validity of the parametric bootstrap for goodness-of-fit testing in semiparametric models*. Annales de l'Institut Henri Poincaré Probabilités et Statistiques, 44(6), 1096--1127. DOI: 10.1214/07-AIHP148.
\item Genest, C., Rémillard, B., \& Beaudoin, D. (2009). *Goodness-of-fit tests for copulas: A review and a power study*. Insurance: Mathematics and Economics, 44(2), 199--213. DOI: 10.1016/j.insmatheco.2007.10.005.
\item Go, Y.-H., \& Lau, W.-Y. (2024). *Terms of trade or market power? Further evidence from dynamic spillovers in return and volatility between Malaysian crude palm oil and foreign exchange markets*. North American Journal of Economics and Finance, 73, 102178. DOI: 10.1016/j.najef.2024.102178.
\item Ignatieva, K., \& Ponomareva, N. (2017). *Commodity currencies and commodity prices: modelling static and time-varying dependence*. Applied Economics, 49(15), 1491--1512. DOI: 10.1080/00036846.2016.1221038.
\item Kojadinovic, I. (2017). *Some copula inference procedures adapted to the presence of ties*. Computational Statistics \& Data Analysis, 112, 24--41. DOI: 10.1016/j.csda.2017.02.006.
\item Koike, T., Kato, S., \& Yoshiba, T. (2025). *Measuring and testing tail equivalence*. Journal of Multivariate Analysis, 209, 105460. DOI: 10.1016/j.jmva.2025.105460.
\item Lo, A. W., \& MacKinlay, A. C. (1990). *An econometric analysis of nonsynchronous trading*. Journal of Econometrics, 45(1--2), 181--211. DOI: 10.1016/0304-4076(90)90098-E.
\item Martens, M., \& Poon, S.-H. (2001). *Returns synchronization and daily correlation dynamics between international stock markets*. Journal of Banking \& Finance, 25(10), 1805--1827. DOI: 10.1016/S0378-4266(00)00159-X.
\item Mensah, P. O., \& Adam, A. M. (2020). *Copula-Based Assessment of Co-Movement and Tail Dependence Structure Among Major Trading Foreign Currencies in Ghana*. Risks, 8(2), 55. DOI: 10.3390/risks8020055.
\item Ning, C., \& Xu, D. (2025). *Extreme comovements and downside/upside risk spillovers between oil prices and exchange rates*. Macroeconomic Dynamics, 29, e53. DOI: 10.1017/S1365100524000543.
\item Qiao, T., \& Han, L. (2023). *COVID-19 and tail risk contagion across commodity futures markets*. Journal of Futures Markets, 43(2), 242--272. DOI: 10.1002/fut.22388.
\item Reboredo, J. C. (2012). *Modelling oil price and exchange rate co-movements*. Journal of Policy Modeling, 34(3), 419--440. DOI: 10.1016/j.jpolmod.2011.10.005.
\item Scholes, M., \& Williams, J. (1977). *Estimating betas from nonsynchronous data*. Journal of Financial Economics, 5(3), 309--327. DOI: 10.1016/0304-405X(77)90041-1.
\item Schotman, P. C., \& Zalewska, A. (2006). *Non-synchronous trading and testing for market integration in Central European emerging markets*. Journal of Empirical Finance, 13(4--5), 462--494. DOI: 10.1016/j.jempfin.2006.04.002.
\item Supper, H., Irresberger, F., \& Weiß, G. (2020). *A comparison of tail dependence estimators*. European Journal of Operational Research, 284(2), 728--742. DOI: 10.1016/j.ejor.2019.12.041.
\item Woode, J. K., Idun, A. A.-A., Kawor, S., Owusu Junior, P., \& Adam, A. M. (2025). *Dynamic Connectedness Between Commodities, Exchange Rates and Equity Markets of Commodity-Dependent Sub-Saharan Africa Countries*. SAGE Open, 15(3). DOI: 10.1177/21582440251347723.
\item Yeboah, S., Fumey, M., Winful, S. A., Otoo, I., \& Owusu, P. (2025). *Asymmetric dependence between commodity prices and selected macroeconomic variables in Ghana*. Scientific African, 28, e02739. DOI: 10.1016/j.sciaf.2025.e02739.
\item Zankawah, M. M., \& Stewart, C. (2020). *Measuring the volatility spill-over effects of crude oil prices on the exchange rate and stock market in Ghana*. Journal of International Trade \& Economic Development, 29(4), 420--439. DOI: 10.1080/09638199.2019.1692895.

\end{itemize}
\textbf{Note:} Bibliographic metadata will be normalized against Crossref and the target journal style before conversion to the final \texttt{references.bib}.

\section{Data and Calendar Construction}
\label{sec:data_and_calendar_construction}

\subsection{Data sources and sample}

The analysis combines daily-source-frequency commodity prices with the Ghanaian cedi exchange rate. The commodity set comprises cocoa, gold, Brent crude oil and West Texas Intermediate (WTI) crude oil. Commodity prices were obtained from Yahoo Finance, while the Cedi series was obtained from the Bank of Ghana. The audited raw files used for the reconstruction have been preserved and hashed as part of the project provenance record.

The analysis begins in January 2015. The two calendar constructions share the same starting period but differ at the end because the underlying sources do not terminate on the same date. The Yahoo commodity data extend beyond the final Bank of Ghana observation, while the last actual Bank of Ghana exchange-rate observation in the audited source is 10 July 2026. The business-day construction therefore continues to 15 July 2026 under its documented limited forward-filling rule, whereas the common-observation construction ends on 10 July 2026.

The Cedi series is expressed so that positive returns correspond to an appreciation of the Cedi and negative returns to a depreciation. The reconstructed sign orientation matches the archived processing exactly, with correlation equal to one and maximum numerical difference below \(2\times10^{-15}\) in the audit comparison.

The resulting return panels are:

\begin{itemize}
\item \textbf{Panel A, business-day alignment:} 3,008 return observations from 5 January 2015 to 15 July 2026.
\item \textbf{Panel B, common-observation alignment:} 2,774 return observations from 5 January 2015 to 10 July 2026.

\end{itemize}
The corresponding marginal-model probability integral transform (PIT) panels begin one observation later because the fitted conditional models require the first return as initialization:

\begin{itemize}
\item \textbf{Panel A PIT:} 3,007 observations from 6 January 2015 to 15 July 2026.
\item \textbf{Panel B PIT:} 2,773 observations from 6 January 2015 to 10 July 2026.

\end{itemize}
The two panels are treated as co-equal representations of the same underlying daily-source information. Neither is labelled the primary, corrected or true panel.

\subsection{Panel A: business-day alignment}

Panel A reproduces the documented business-day construction used in the original analysis. A business-day calendar is imposed across the five series and missing prices are carried forward for at most three business days before log returns are calculated. The purpose of this construction is to retain a nominal business-day frequency while avoiding unrestricted interpolation across long gaps.

The rule has observable consequences. Because a carried-forward price does not change, it generates a zero return for that asset on the affected business day. In the audited Panel A return file, the fraction of exact zero returns is approximately 4.42\% for cocoa, 4.06\% for gold, 4.06\% for Brent, 4.09\% for WTI and 22.27\% for the Cedi. The much larger zero share for the Cedi reflects both its underlying quote process and the interaction between the quote calendar and the business-day alignment rule. For this reason, the Cedi margin is treated separately in the diagnostic discussion rather than assumed to behave like the commodity margins.

Panel A contains 3,008 return rows. Its reconstructed return file has SHA256 checksum

\texttt{3e9769f42ec7925aacd951277ddae6ff529822408815b01210de87659221341f}.

The audited Panel A PIT file has SHA256 checksum

\texttt{be42455c12bf2f4a9c7ef33494930ad396f569192f571386d099046a388119d1}.

These hashes are used as fixed provenance checks for later analysis scripts.

The end of the Panel A sample illustrates the implications of this construction. The audited Yahoo data extend beyond the last actual Bank of Ghana observation, while the Bank of Ghana raw series ends on 10 July 2026. Under the documented three-business-day carry-forward rule, the Cedi price is therefore carried into 13, 14 and 15 July, allowing the aligned panel to end on 15 July 2026. These final dates are not additional Bank of Ghana observations; they arise from the calendar rule.

\subsection{Panel B: common-observation alignment}

Panel B is constructed without forward-filled return observations. Instead, prices are restricted to dates on which all five series have actual source observations, and returns are calculated between consecutive common observation dates.

For asset \(i\), let \(t^{-}\) denote the previous date on which all five series are jointly observed. The return is

\[
r_{i,t}=100\left[\log P_{i,t}-\log P_{i,t^{-}}\right].
\]

For the Cedi, the sign is reversed to maintain the interpretation that positive returns correspond to appreciation.

This construction produces synchronized close-to-close returns on common observation dates, but it should not be described as perfect intraday synchronization or as a fixed 24-hour daily return series. The interval between consecutive common dates varies with weekends, holidays and source-specific missing observations. Among the 2,774 Panel B returns, 2,118 intervals span one calendar day, 53 span two days, 461 span three days, 115 span four days, 26 span five days and one spans six days. The median interval is one day and the maximum is six days.

Before the final return calculation, 2,776 dates are common to the five actual-source price series. One synchronized endpoint, 21 April 2020, is invalid because the preceding common-price interval crosses the non-positive WTI observation of 20 April 2020. After removing the affected endpoint, the final Panel B contains 2,774 returns from 5 January 2015 to 10 July 2026.

The Panel B return file has SHA256 checksum

\texttt{e6312d6163aeb7c9fb7ce722ebe46be1898f23855d5a87395830e4d1e82bf3b3}.

Its PIT file has SHA256 checksum

\texttt{9b322bda7c4e63534da14d1a7e360e758762e98cb39101cf329cdceb5afc216f}.

As with Panel A, these hashes are used as fixed reproducibility checks.

\subsection{Negative WTI price episode}

WTI presents a special data issue on 20 April 2020, when the audited raw price is \(-37.630001\). A standard log return is undefined at a non-positive price, so this observation cannot enter the return calculation in the same way as an ordinary positive price.

The two calendar constructions handle the episode differently because their alignment rules differ. In Panel A, the processed positive-price sequence around the episode is:

\begin{itemize}
\item 17 April 2020: 18.27, return \(-8.3951\%\);
\item 20 April 2020: 18.27, return \(0\%\);
\item 21 April 2020: 10.01, return \(-60.1676\%\);
\item 22 April 2020: 13.78, return \(31.9634\%\).

\end{itemize}
Thus the negative raw WTI observation is not used directly in the Panel A log-return calculation. The resulting loss is shifted to the next positive-price observation rather than represented as a log return on 20 April itself.

Panel B uses only common actual-source price dates and requires a valid positive-price interval. The 20 April non-positive observation makes the synchronized interval ending on 21 April invalid, so both 20 and 21 April are absent from the final Panel B return series. The next retained WTI return is 22 April 2020, \(31.9634\%\).

The paper reports this episode explicitly because it is economically exceptional and because different preprocessing rules can relocate or remove extreme observations. No attempt is made to treat the negative settlement as if it were a standard positive log-price observation.

\subsection{Zero returns and small movements}

The two panels differ markedly in the frequency of exact zeros. In Panel B, exact zero-return shares are much lower for the commodities and lower for the Cedi than under business-day alignment:

\begin{table}[htbp]
\centering
\small
\begin{tabularx}{\textwidth}{lXX}
\toprule
\textbf{Series} & \textbf{Exact zero returns, Panel A} & \textbf{Exact zero returns, Panel B} \\
\midrule
Cocoa & 4.42\% & 0.65\% \\
Gold & 4.06\% & 0.40\% \\
Brent & 4.06\% & 0.40\% \\
WTI & 4.09\% & 0.43\% \\
Cedi & 22.27\% & 16.51\% \\
\bottomrule
\end{tabularx}
\end{table}

Panel B also retains a high concentration of small Cedi movements. Approximately 47.19\% of non-zero Cedi returns have absolute magnitude at most 0.05\%, 61.14\% are at most 0.10\%, and 89.98\% are at most 0.50\%. The zero-heavy and highly concentrated character of the Cedi return distribution is therefore not created solely by business-day forward filling, although the business-day construction increases the exact-zero frequency.

This feature matters for the marginal modeling stage. The Cedi's selected continuous marginal specification fails PIT adequacy diagnostics in both calendar panels, while the selected commodity margins pass their stated adequacy gates. Consequently, later Cedi-related copula and extreme-value quantities are presented as conditional diagnostics and interpreted more cautiously than the commodity-only results.

\subsection{Descriptive comparison of the two panels}

Although the calendar constructions differ, the resulting return series remain strongly related on shared dates. Pearson correlations between corresponding Panel A and Panel B series are approximately 0.976 to 0.986, while Spearman correlations are approximately 0.975 to 0.996. The panels should therefore not be understood as two unrelated datasets. They are two closely related representations of the same market information that differ in how they handle non-common trading dates and missing observations.

The stress-window sample counts also change. Panel A contains 87 observations in the prespecified COVID-19 window and 262 observations in 2024. Panel B contains 80 COVID-19 observations and 242 observations in 2024. The combined Panel B stress sample therefore contains 322 observations. These differences are part of the design rather than nuisances to be hidden, because the paper asks whether inferential conclusions depend on the calendar construction that generates them.

\subsection{Why both panels are retained}

Neither construction dominates the other on every criterion. Panel A preserves a regular business-day index and closely reproduces the documented historical pipeline, but forward filling can generate artificial zero returns and can extend an aligned series beyond the final actual observation of one source. Panel B avoids forward-filled return observations and ensures that each multivariate row uses the same pair of common observation dates, but its returns span irregular one- to six-day calendar intervals and therefore cannot be interpreted as fixed-horizon 24-hour returns.

For this reason, the paper does not designate one panel as primary and the other as a robustness check. The two are co-equal specifications of the observation process. A substantive conclusion is called \textbf{calendar-robust} when the same inferential conclusion is obtained under both constructions. A conclusion is called \textbf{calendar-sensitive} when the two constructions lead to different inferential classifications. Findings that arise only as diagnostics of one construction are labelled \textbf{panel-specific}.

This classification is determined by the analysis protocol rather than by which panel produces the more convenient result. It is used consistently in the copula adequacy and stress-inference sections that follow.

\subsection{Reproducibility and provenance}

The reconstruction preserves raw-source hashes, processed-return hashes and PIT hashes so that later analysis can be tied to fixed inputs. The data-audit stage also records special-case transformations such as the WTI negative-price episode, the Cedi sign orientation and the dates added to Panel A through limited forward filling.

The paper intentionally avoids making claims about unverified details of the upstream commodity contract construction that are not established by the archived raw data and scripts. In particular, the manuscript does not assert an unverified continuous-futures roll or back-adjustment rule. The publication record is restricted to processing steps that were directly audited and reproduced.

\paragraph{Table 1. Calendar construction summary}

\begin{table}[htbp]
\centering
\small
\begin{tabularx}{\textwidth}{lXX}
\toprule
\textbf{Feature} & \textbf{Panel A: business-day alignment} & \textbf{Panel B: common-observation alignment} \\
\midrule
Return observations & 3,008 & 2,774 \\
Return period & 5 Jan 2015--15 Jul 2026 & 5 Jan 2015--10 Jul 2026 \\
PIT observations & 3,007 & 2,773 \\
PIT period & 6 Jan 2015--15 Jul 2026 & 6 Jan 2015--10 Jul 2026 \\
Price alignment & Business-day grid & Common actual-source dates \\
Missing-price rule & Forward fill, maximum 3 business days & No forward-filled return observations \\
Return interval & Nominal business-day & Previous common observation date \\
Calendar-day span & Normally one business-day step, with carried prices possible & 1--6 calendar days, median 1 \\
Final actual BoG date & 10 Jul 2026 & 10 Jul 2026 \\
Final panel date & 15 Jul 2026 & 10 Jul 2026 \\
WTI 20 Apr 2020 & Negative raw price excluded from direct log return; processed price held at 18.27 & Non-positive interval causes affected common endpoint to be dropped \\
Exact Cedi zero-return share & 22.27\% & 16.51\% \\
Role in paper & Co-equal panel & Co-equal panel \\
\bottomrule
\end{tabularx}
\end{table}

\textbf{Note:} "Common-observation alignment" means returns between consecutive dates on which all five source series are observed. It does not imply identical intraday timestamps.

\section{Methodology}
\label{sec:methodology}

\subsection{Marginal filtering and the validity gate}

Each return series is filtered before dependence estimation using a fixed four-model ladder: AR(1)-GJR-GARCH(1,1)-t, AR(2)-GJR-GARCH(1,1)-t, AR(1)-EGARCH(1,1)-t, and AR(1)-GJR-GARCH(1,1)-skew-t. For the GJR-GARCH candidates,

\[
r_t=\mu+\sum_{j=1}^{p}\phi_j r_{t-j}+\varepsilon_t,\qquad
\varepsilon_t=\sigma_t z_t,
\]

\[
\sigma_t^2=\omega+\left(\alpha+\gamma I\{\varepsilon_{t-1}<0\}\right)\varepsilon_{t-1}^2+\beta\sigma_{t-1}^2.
\]

The ladder is fixed rather than expanded after seeing the data.

A candidate is eligible only if the optimizer reports \texttt{convergence\_flag = 0} and the fitted log likelihood, parameters, standardized residuals, conditional volatilities and PIT values are finite. Non-converged candidates do not enter either the adequacy pool or the AIC fallback.

For each valid fit, standardized residuals are \(\hat z_t=\hat\varepsilon_t/\hat\sigma_t\). Adequacy is assessed at the 5\% level using a Ljung-Box test on standardized residuals, a Ljung-Box test on squared standardized residuals, and a Kolmogorov-Smirnov test of PIT values against Uniform(0,1). The Ljung-Box diagnostics use 10 lags. A candidate is adequate only when all three tests fail to reject.

The first valid adequate candidate in the fixed ladder is selected. If none passes, the valid converged candidate with the lowest AIC is retained only as an \textbf{inadequate reference model}. This distinction is preserved in downstream interpretation.

\subsection{PIT values and rank-based pseudo-observations}

For a selected margin,

\[
\tilde u_t=\hat F(\hat z_t),
\]

where \(\hat F\) is the fitted innovation CDF. PIT values are clipped to \([10^{-6},1-10^{-6}]\) for numerical stability. Copula estimation then uses rank-based pseudo-observations,

\[
u_t=\frac{\operatorname{rank}(\tilde u_t)}{n+1}.
\]

The clipping rule was separately audited. In Panel A it creates one three-way Cedi floor tie; recovering the underlying order changes no prespecified tail-set membership and only negligibly changes observed Cedi-pair GOF statistics.

\subsection{Extreme-value diagnostics}

Extreme-value calculations are applied to the loss transform of standardized residuals,

\[
L_t=-\hat z_t.
\]

A peaks-over-threshold generalized Pareto model is fitted above the empirical 90th percentile of \(L_t\). For exceedance \(Y=L-u\),

\[
G_{\xi,\beta}(y)=1-\left(1+\xi\frac{y}{\beta}\right)^{-1/\xi},
\]

on its admissible support. The analysis records the threshold, shape \(\xi\), scale \(\beta\), number of exceedances, 99\% VaR, 99\% expected shortfall and a Hill diagnostic using \(k=\max\{25,\lfloor0.05n\rfloor\}\). The main threshold remains fixed at 0.90; Cedi threshold sensitivity is treated as a diagnostic, not as an alternative headline specification.

\subsection{Pairwise copula families}

With five series there are ten unordered pairs. Each pair is estimated under eight static copula families: Gaussian, Student-t, Clayton, Gumbel, Frank, Joe, survival Clayton and survival Gumbel.

AIC is retained only as a relative-fit descriptor. Absolute adequacy is assessed separately for every pair-family cell. A dedicated implementation audit confirmed that the Panel A and Panel B GOF branches use equivalent family estimators when supplied with the same pseudo-observations.

\subsection{Absolute copula goodness of fit}

Absolute GOF uses a Rosenblatt-transform Cramér-von Mises statistic with parametric bootstrap. For fitted family \(C_{\hat\theta}\),

\[
e_{1i}=u_i,\qquad e_{2i}=h_{\hat\theta}(v_i\mid u_i).
\]

The transformed empirical CDF is evaluated on a 50 by 50 grid,

\[
g_j=\frac{j}{51},\qquad j=1,\ldots,50,
\]

and

\[
S_n=
\sum_{j=1}^{50}\sum_{k=1}^{50}
\left[\hat D_n(g_j,g_k)-g_jg_k\right]^2.
\]

For each bootstrap replicate, a sample of size \(n\) is simulated from the fitted family, the simulated margins are re-ranked, the family is re-estimated, and the statistic is recomputed. The Monte Carlo p-value is

\[
\hat p=\frac{1+\sum_{b=1}^{B}I\{S_n^{(b)}\ge S_n\}}{B+1}.
\]

All 80 cells are first evaluated with \(B=2000\). Cells with initial \(0.03\le\hat p\le0.07\) are automatically rerun with \(B=5000\), and the higher-resolution result replaces the base value. The implemented procedure is described as a \textbf{Rosenblatt-transform Cramér-von Mises statistic with parametric bootstrap}. It is not claimed to be an exact reproduction of every statistic in \citet{genest2009}.

A rejection at 5\% means the tested family is inadequate under this procedure. Non-rejection is not evidence that the family is the true model.

\subsection{Finite-threshold tail concentration}

For pseudo-observations \(u_i,v_i\),

\[
\hat\lambda_L(q)=\frac{1}{nq}\sum_{i=1}^{n}I(u_i\le q,v_i\le q),
\]

\[
\hat\lambda_U(q)=\frac{1}{nq}\sum_{i=1}^{n}I(u_i\ge1-q,v_i\ge1-q),
\]

with \(q\in\{0.025,0.05,0.10\}\).

These are finite-threshold concentration measures, not direct estimates of the asymptotic tail-dependence coefficient as \(q\to0\).

\subsection{Stress definitions and calm baseline}

Three stress families are evaluated: combined COVID-19 plus 2024, COVID-19 alone from 1 March to 30 June 2020, and calendar year 2024 alone. All three use the same calm baseline, which excludes both stress windows.

Within every calm and stress subsample, PIT values are re-ranked before tail concentration is calculated. For each pair, tail and threshold,

\[
d=\hat\lambda_{stress}(q)-\hat\lambda_{calm}(q).
\]

Each family contains 10 pairs x 3 thresholds x 2 tails = 60 comparisons.

\subsection{Moving-block bootstrap}

Inference for \(d\) uses a moving-block bootstrap with block length 20 and \(B=15000\) replications per cell. Calm and stress sequences are resampled separately. Within each bootstrap replicate, both margins are re-ranked before the finite-tail statistic is recalculated.

If \(d^{(b)}\) is the bootstrap difference, the two-sided centered p-value is

\[
\hat p=
\frac{1+\sum_{b=1}^{B}I\{|d^{(b)}-d|\ge |d|\}}{B+1}.
\]

The bootstrap output also records the empirical 2.5\% and 97.5\% quantiles of the difference distribution. Stable SHA256-derived seeds are used for reproducibility.

\subsection{Multiplicity control}

Multiplicity is controlled separately within each 60-cell stress family. Bonferroni control uses

\[
\alpha_{Bonf}=0.05/60\approx0.0008333,
\]

and the Benjamini-Hochberg procedure controls the false discovery rate at 5\%.

The combined, COVID-only and 2024-only analyses are therefore three separate correction families, not one 180-test family. Nominal p-values below 0.05 remain descriptive unless they survive the relevant correction.

\subsection{Cross-panel interpretation}

The two calendar panels are co-equal. A result is \textbf{calendar-robust} when the same inferential classification is obtained in both panels under the same rule. For a cell-level stress discovery, this requires survival of the same multiplicity procedure in both panels and the same direction of the estimated difference.

A result is \textbf{calendar-sensitive} when corrected status changes across panels or the direction changes in a way that alters interpretation. A result is \textbf{panel-specific} when it is reported only as a diagnostic of one construction.

Nominal-only differences are not promoted to robust discoveries, and GOF non-rejection is not interpreted as proof of correctness.

\subsection{Reproducibility controls}

The publication analysis uses fixed input hashes, deterministic seeds, checkpointed long-running computations and non-destructive output directories. Panel-specific PIT files are hash-checked before downstream analysis. Historical exploratory artifacts are retained for provenance but do not override results reproduced from the frozen inputs under the audited implementation.

\section{Results}
\label{sec:results}

\subsection{Marginal-model adequacy across the two calendar panels}

The marginal stage produces a mixed but internally consistent result across the two calendar constructions. The selected commodity margins are adequate under the stated three-test gate in both panels, although the selected specification is not always the same. The Cedi is the exception: no tested candidate passes the adequacy gate in either panel.

\begin{table}[htbp]
\centering
\small
\begin{tabularx}{\textwidth}{lXXXX}
\toprule
\textbf{Series} & \textbf{Panel A: business-day alignment} & \textbf{Status} & \textbf{Panel B: common-observation alignment} & \textbf{Status} \\
\midrule
Cocoa & AR(1)-GJR-GARCH-t & Adequate & AR(1)-GJR-GARCH-t & Adequate \\
Gold & AR(1)-GJR-GARCH-t & Adequate & AR(1)-GJR-GARCH-t & Adequate \\
Brent & AR(1)-GJR-GARCH-skew-t & Adequate & AR(1)-GJR-GARCH-t & Adequate \\
WTI & AR(1)-EGARCH-t & Adequate & AR(1)-GJR-GARCH-t & Adequate \\
Cedi & AR(1)-GJR-GARCH-t & Inadequate reference & AR(1)-EGARCH-t & Inadequate reference \\
\bottomrule
\end{tabularx}
\end{table}

The Panel A hard-gate audit checked all 20 candidate fits. One candidate, the Cedi AR(1)-EGARCH-t specification, was invalid because the optimizer did not converge. Imposing convergence and finite-fit validity as hard requirements did not change any of the previously selected Panel A models. The Panel B selection differs for Brent, WTI and the Cedi, but the higher-level adequacy pattern is unchanged: the four commodity margins have an adequate selected model, while the Cedi does not.

This result is important for the interpretation that follows. Commodity-only dependence results are supported by margins that pass the stated gate in both panels. Cedi-related copula and extreme-value quantities are conditional on an inadequate continuous marginal reference under both calendar constructions.

\subsection{Static copula adequacy: robust failures, robust non-rejections and calendar sensitivity}

The clearest cross-panel difference appears in the absolute copula goodness-of-fit results. Table 2 reports the number of rejected families, out of the eight tested, for each pair.

\begin{table}[htbp]
\centering
\small
\begin{tabularx}{\textwidth}{lXXX}
\toprule
\textbf{Pair} & \textbf{Panel A rejected} & \textbf{Panel B rejected} & \textbf{Cross-panel classification} \\
\midrule
Cocoa-Gold & 8/8 & 2/8 & Calendar-sensitive \\
Cocoa-Brent & 8/8 & 3/8 & Calendar-sensitive \\
Cocoa-WTI & 8/8 & 3/8 & Calendar-sensitive \\
Cocoa-Cedi & 0/8 & 0/8 & Calendar-robust non-rejection \\
Gold-Brent & 8/8 & 8/8 & Calendar-robust rejection \\
Gold-WTI & 8/8 & 8/8 & Calendar-robust rejection \\
Gold-Cedi & 0/8 & 0/8 & Calendar-robust non-rejection \\
Brent-WTI & 8/8 & 8/8 & Calendar-robust rejection \\
Brent-Cedi & 0/8 & 0/8 & Calendar-robust non-rejection \\
WTI-Cedi & 0/8 & 0/8 & Calendar-robust non-rejection \\
\bottomrule
\end{tabularx}
\end{table}

Three commodity-commodity pairs are therefore robustly inconsistent with every member of the tested eight-family static copula set: Gold-Brent, Gold-WTI and Brent-WTI. This conclusion survives both calendar constructions.

The four commodity-Cedi pairs show the opposite pattern. None of the eight tested static families is rejected for Cocoa-Cedi, Gold-Cedi, Brent-Cedi or WTI-Cedi in either panel. This is also calendar-robust, but its interpretation is deliberately limited. Non-rejection does not establish that all eight families are equally plausible or that any one of them is the true dependence model. It only states that the implemented GOF procedure does not reject them at the 5\% level.

The cocoa-related commodity pairs are where the calendar construction matters most. Panel A rejects all eight families for Cocoa-Gold, Cocoa-Brent and Cocoa-WTI. Panel B rejects only two of eight for Cocoa-Gold and three of eight for Cocoa-Brent and Cocoa-WTI. The broad claim that every commodity-commodity pair rejects every tested static copula is therefore not robust. The evidence supports a narrower statement: complete eight-family rejection is robust for Gold-Brent, Gold-WTI and Brent-WTI, while the three cocoa-related commodity pairs are calendar-sensitive.

One Panel A Cocoa-Cedi Frank result lies close to the 5\% decision boundary after the higher-resolution bootstrap. The later PIT-clipping audit shows that its observed GOF statistic is unchanged when the artificial three-way Cedi floor tie is de-clipped. The borderline classification is therefore a Monte Carlo-precision issue rather than a consequence of the PIT-clipping artifact.

\subsection{Stress-period finite-tail concentration}

The stress analysis produces a different kind of robustness result. Across the combined COVID-19 plus 2024 definition, COVID-19 alone and 2024 alone, neither panel produces a multiplicity-adjusted discovery under Bonferroni or Benjamini-Hochberg correction.

\begin{table}[htbp]
\centering
\small
\begin{tabularx}{\textwidth}{lXXXXXXX}
\toprule
\textbf{Stress definition} & \textbf{Panel A nominal p<0.05} & \textbf{Panel B nominal p<0.05} & \textbf{Panel A Bonferroni} & \textbf{Panel B Bonferroni} & \textbf{Panel A BH} & \textbf{Panel B BH} & \textbf{Classification} \\
\midrule
Combined COVID + 2024 & 6/60 & 5/60 & 0 & 0 & 0 & 0 & Calendar-robust \\
COVID-19 only & 11/60 & 11/60 & 0 & 0 & 0 & 0 & Calendar-robust \\
2024 only & 6/60 & 7/60 & 0 & 0 & 0 & 0 & Calendar-robust \\
\bottomrule
\end{tabularx}
\end{table}

The nominal counts show that several individual cells fall below 0.05 before correction, and the exact set is not identical across calendars. Those nominal results are descriptive rather than headline discoveries. The inferential conclusion is determined within each prespecified 60-test family.

Under that rule, the result is the same in all six panel-by-stress combinations: zero Bonferroni discoveries and zero Benjamini-Hochberg discoveries. The absence of multiplicity-adjusted stress findings is therefore calendar-robust under the tested stress definitions, thresholds and moving-block-bootstrap design.

This result should not be restated as evidence that tail dependence is constant or that crises do not affect commodity or exchange-rate relationships. The statistical result is narrower: none of the 60 finite-threshold stress-versus-calm comparisons in any of the three stress families survives either multiplicity procedure in either panel.

\subsection{Cedi-related stress evidence and the prespecified follow-up trigger}

The absence of corrected stress findings also applies to commodity-Cedi cells in Panel B. The prespecified rule required a full synchronized seven-treatment Cedi robustness rerun only if at least one Panel B commodity-Cedi stress cell survived Bonferroni or Benjamini-Hochberg correction.

That trigger was not activated under the combined, COVID-only or 2024-only stress definition. Accordingly, the existing seven-treatment Panel A analysis remains supporting robustness evidence, but it is not promoted to a symmetric two-panel result. No additional synchronized Cedi-treatment search is needed to sustain the current headline stress conclusion.

\subsection{Extreme-value diagnostics}

The EVT results are interpreted separately from the copula stress tests. In Panel B, the 90th-percentile POT/GPD fit for the Cedi gives approximately

\[
\hat\xi=0.693,\qquad
\hat\beta=0.329,
\]

with

\[
\operatorname{VaR}_{0.99}\approx2.332
\]

and

\[
\operatorname{ES}_{0.99}\approx7.616
\]

in standardized-residual units. Across the Cedi threshold-sensitivity range from 0.85 to 0.975, the fitted shape parameter remains positive and ranges approximately from 0.63 to 0.79.

These values indicate a heavy fitted loss tail under the retained Panel B reference model, but the Cedi margin fails the adequacy gate. The estimates are therefore conditional diagnostics. They should not be presented as model-free evidence that the Cedi has a particular structural tail index, and they should not be substituted for Panel A quantities without identifying the calendar panel.

The same caution applies to the very large degrees-of-freedom estimate in the Panel A Cocoa-Cedi Student-t copula. The PIT-tie audit shows that the estimate is numerically weakly identified near the Gaussian limit and that the associated GOF statistic is essentially unchanged by recovering the unclipped rank ordering. The large numerical movement in the degrees-of-freedom parameter is not a substantive change in dependence.

\subsection{Results by research question}

The three research questions can now be answered directly from the cross-panel evidence.

\textbf{RQ1: Do finite-threshold lower- and upper-tail concentration measures differ between calm and stress periods, and are those conclusions robust to calendar alignment?}  
At the nominal 5\% level, several cells differ from calm in each stress family. After the prespecified Bonferroni and Benjamini-Hochberg corrections, no cell survives in either panel for the combined, COVID-only or 2024-only stress definition. The corrected stress conclusion is therefore calendar-robust.

\textbf{RQ2: Can any of the eight tested static bivariate copula families adequately represent the dependence structure, and which conclusions are robust across calendars?}  
For Gold-Brent, Gold-WTI and Brent-WTI, all eight tested families are rejected in both panels. For all four commodity-Cedi pairs, none of the eight families is rejected in either panel. Cocoa-Gold, Cocoa-Brent and Cocoa-WTI are calendar-sensitive because complete rejection in Panel A weakens to partial rejection in Panel B.

\textbf{RQ3: How sensitive are Cedi-related dependence and extreme-value diagnostics to the Cedi's zero-heavy distribution, marginal inadequacy and calendar construction?}  
The Cedi remains inadequately modeled by the tested continuous marginal candidates under both panels. Its exact selected reference model changes across calendars, and Panel B EVT estimates remain strongly heavy-tailed, but those EVT quantities are conditional on marginal misspecification. The Panel A PIT-clipping audit shows that the identified three-value floor tie does not alter the prespecified stress-tail memberships and has negligible effect on observed Cedi-pair GOF statistics. The main Cedi-related limitation is therefore the inadequacy of the tested marginal models, not the small clipping artifact.

\subsection{Empirical synthesis}

Taken together, the results separate three different kinds of evidence.

First, there is \textbf{calendar-robust copula inadequacy} for Gold-Brent, Gold-WTI and Brent-WTI within the eight tested static families.

Second, there is \textbf{calendar-sensitive copula adequacy} for Cocoa-Gold, Cocoa-Brent and Cocoa-WTI. These pairs show that calendar construction can materially alter absolute-model conclusions even when the underlying return series remain highly correlated across panels.

Third, there is a \textbf{calendar-robust absence of multiplicity-adjusted stress discoveries} across all three stress definitions. The nominal p-value patterns vary, but no corrected stress finding survives in either panel.

The Cedi results require a separate qualification: the copula non-rejections are cross-panel robust, but the Cedi marginal remains inadequate in both constructions. Cedi-related dependence and EVT results are therefore retained as conditional evidence rather than used to support stronger structural claims.

\section{Discussion}
\label{sec:discussion}

\subsection{Stress-period evidence: nominal signals without multiplicity-adjusted discoveries}

The stress results are more informative when read as an inferential result than as a simple count of p-values. Across the combined COVID-19 plus 2024 definition, COVID-19 alone and 2024 alone, both calendar panels contain several nominal \(p<0.05\) cells. None survives either Bonferroni or Benjamini-Hochberg correction. The important result is therefore not that the estimated tail-concentration differences are all zero. They are not. It is that, under the prespecified 60-test families and the moving-block-bootstrap design, the observed differences do not provide multiplicity-adjusted evidence strong enough to support individual stress effects.

This finding is narrower than some of the crisis literature, but it is not necessarily inconsistent with it. \citet{qiao2023} report stronger lower- and upper-tail contagion across a broad network of commodity futures after the onset of COVID-19, while other dynamic and quantile-based studies document stronger cross-market connectedness during periods of distress. Those studies examine different asset sets, dependence objects, model structures and testing procedures. The present study asks a more specific question: whether finite-threshold concentration for ten Ghana-relevant bivariate pairs differs between the prespecified stress and calm regimes strongly enough to survive correction across 60 simultaneous comparisons.

That distinction matters. A crisis can alter market dependence without producing a corrected discovery in every pair, tail and threshold. Conversely, a collection of nominal \(p<0.05\) values can arise in a large comparison set without establishing a robust crisis effect. The absence of Bonferroni and BH discoveries in both panels therefore supports a cautious conclusion: within the tested thresholds, windows and bootstrap design, the data do not provide multiplicity-adjusted evidence for stress-period changes in finite-tail concentration.

The agreement between the two calendar panels strengthens that conclusion. The exact nominal cells differ, but the family-level inferential outcome is the same in all three stress definitions. In this sense, the stress result is more robust than the individual nominal p-values that generate it.

\subsection{Why calendar construction matters most for the cocoa pairs}

The copula GOF results show that calendar construction can matter even when two return panels are highly correlated. Gold-Brent, Gold-WTI and Brent-WTI reject all eight tested static families in both panels, whereas Cocoa-Gold, Cocoa-Brent and Cocoa-WTI reject all eight families in the business-day panel but only a subset in the common-observation panel.

This is consistent with the broader nonsynchronous-trading literature. \citet{scholes1977} and \citet{lo1990} show that the observation process can affect return moments and cross-dependence. \citet{martens2001} demonstrate that close-to-close returns measured across different market hours can distort international correlation estimates, and \citet{schotman2006} show that moving to lower frequency does not automatically remove the problem. \citet{dimitriou2025} provide recent evidence that synchronization choices can alter spillover conclusions during crisis periods.

The present results extend that concern from correlation and spillover measurement to absolute copula adequacy. The two panels remain strongly correlated on shared dates, yet the rejection pattern for all three cocoa-related commodity pairs changes materially. This means that a high cross-panel correlation is not sufficient to establish that higher-order dependence features are invariant to calendar construction.

The result does not identify one calendar as correct. Panel A preserves a regular business-day structure but introduces zero returns through limited carry-forward. Panel B removes forward-filled return observations but allows variable one- to six-day intervals between common observations. Both are defensible representations of imperfectly synchronized source data. The empirical contribution is therefore the identification of which conclusions survive both representations and which do not.

For cocoa-related dependence in particular, the correct interpretation is calendar sensitivity, not model failure under one preferred dataset. Any substantive interpretation of Cocoa-Gold, Cocoa-Brent or Cocoa-WTI based on a single alignment rule would hide an important source of specification uncertainty.

\subsection{Static copula inadequacy as model risk}

The robust rejection of all eight tested static families for Gold-Brent, Gold-WTI and Brent-WTI is a different result from the stress analysis. It does not say that the pairs are independent, nor does it say that "copulas fail." It says that none of the eight tested static bivariate families provides an adequate representation under the implemented Rosenblatt-transform GOF procedure in either calendar panel.

This distinction between relative fit and absolute adequacy is central. An AIC winner must exist whenever several estimable candidate models are compared, even when every candidate is misspecified. The GOF results show why selecting the lowest-AIC family is not sufficient evidence for a satisfactory dependence model. For the three robustly rejected commodity pairs, model choice within the eight-family set is therefore a choice among relatively better fitting but absolutely inadequate static specifications.

The result is consistent with literature showing that commodity dependence can be time-varying, asymmetric and regime-dependent. \citet{ignatieva2017}, \citet{ning2025}, and related commodity-currency studies document changing dependence across market states. \citet{supper2020} further show that static tail-dependence estimators can be misleading when extreme dependence varies over time. The present GOF results do not identify which alternative class is correct, but they provide a concrete reason to consider richer models such as dynamic copulas, regime-switching structures, vine constructions or nonparametric dependence estimators in future work.

The calendar-sensitive cocoa results add another layer to this model-risk interpretation. A model can appear inadequate under one defensible observation scheme and not under another. Model adequacy is therefore partly a property of the processed statistical object being fitted, not only of the named market pair.

\subsection{Commodity-Cedi non-rejection requires a different interpretation}

The commodity-Cedi pairs produce 0/8 GOF rejections in both panels. It would be incorrect to read this as evidence that all eight families describe commodity-Cedi dependence well. A non-rejection only means that the implemented test does not detect sufficient lack of fit at the chosen level.

There are at least two reasons for caution. First, the empirical dependence signal may be weak enough that very different copula shapes are difficult to distinguish statistically. Second, and more importantly for this study, the Cedi marginal fails the stated adequacy gate in both calendar constructions. \citet{fantazzini2009} shows that marginal misspecification can distort copula and risk estimates. The commodity-Cedi GOF results therefore sit downstream of a known marginal-model limitation.

The Cedi's zero-heavy return process is not merely an artifact of Panel A forward filling. The exact-zero share falls from 22.27\% in Panel A to 16.51\% in Panel B, but it remains large in the common-observation construction. This is consistent with the conclusion that the difficulty is structural to the observed daily Cedi series and the continuous marginal family being imposed, rather than being created solely by the business-day calendar rule.

The clipping audit helps separate a minor numerical issue from the larger modeling problem. The three Panel A Cedi PIT values clipped to \(10^{-6}\) produce one artificial tie, but recovering their underlying order changes no prespecified tail membership and has negligible effect on observed Cedi-pair GOF statistics. The main limitation is therefore not the clipping rule. It is that none of the tested continuous Cedi marginal candidates provides an adequate PIT representation in either panel.

This distinction also governs the EVT interpretation. The positive and relatively large Panel B Cedi GPD shape estimate is evidence of a heavy fitted loss tail conditional on the retained reference model. It is not a model-free estimate of the Cedi's structural tail index. A discrete-continuous or otherwise more flexible marginal model may alter that estimate.

\subsection{Relation to the Ghana literature}

The findings fit more naturally with the Ghana literature when frequency and estimand are kept explicit. \citet{archer2022}, \citet{barson2023}, \citet{yeboah2025} and related studies find asymmetric, nonlinear or time-frequency relationships between commodity prices and the Cedi or Ghanaian macro-financial variables. Those studies do not imply that every daily tail-concentration contrast should be statistically significant after correction.

\citet{barson2023}, for example, emphasize time-frequency comovement, while \citet{atipaga2025} find that Cedi relationships with global uncertainty measures vary across horizons. \citet{woode2025} likewise report time- and frequency-varying connectedness in commodity-dependent sub-Saharan African markets. These findings make it plausible that economically meaningful relationships may appear more clearly at medium or longer horizons than in the daily finite-tail contrasts used here.

The present results should therefore not be framed as contradicting the Ghana literature. They refine the question. The paper finds that, at the daily-source frequency and under the prespecified finite-tail stress design, no individual stress-calm difference survives multiplicity correction across either calendar construction. At the same time, the copula GOF analysis shows substantial dependence-model inadequacy for several commodity pairs. Weak corrected stress evidence and strong model-adequacy concerns can coexist.

\citet{mensah2020} is especially relevant because it already demonstrates that copula-based tail analysis is feasible for Ghanaian exchange rates. The present study differs by combining commodity-Cedi pairs with explicit calendar sensitivity, hard marginal adequacy gates, eight-family absolute GOF testing and multiplicity-controlled stress inference. The contribution is therefore not the introduction of copulas to Ghana, but the identification of which dependence conclusions remain stable after these additional layers of scrutiny.

\subsection{Implications for empirical risk analysis}

Three practical implications follow from the results.

First, crisis labels should not substitute for formal stress inference. COVID-19 and the 2024 episode are economically important periods, but the presence of a named crisis does not guarantee a statistically robust increase in every tail-dependence measure. Risk analysis is better served by prespecified windows, explicit estimands and multiplicity control than by selecting visually striking crisis-period estimates after the fact.

Second, model selection and model adequacy should be reported separately. For Gold-Brent, Gold-WTI and Brent-WTI, every tested static family is rejected in both panels. Reporting only an AIC-best copula for those pairs would create a false sense of precision. A relative winner can still be an inadequate model.

Third, calendar construction should be treated as part of the empirical specification. The cocoa-pair results show that dependence-model conclusions can change even when the corresponding return series remain highly correlated across alternative panels. For applications involving nonsynchronous commodity and exchange-rate data, documenting the alignment rule is therefore part of documenting the statistical model.

These are methodological implications rather than policy prescriptions. The paper does not estimate fiscal losses, reserve effects, hedging effectiveness or causal pass-through from commodity prices to the Cedi. Those questions require separate designs and, in some cases, lower-frequency macroeconomic data.

\subsection{Limitations revealed by the reconstruction}

The reconstruction makes several limitations explicit.

The first is the Cedi marginal. No tested continuous specification passes the adequacy gate under either panel. This limits the strength of commodity-Cedi copula and EVT claims.

The second is the use of static bivariate copulas. The robust GOF rejections for three commodity pairs suggest that the tested static families are too restrictive for some relationships. Dynamic, regime-switching, vine or nonparametric models are natural extensions, but the present results do not establish which of these would be adequate.

The third is calendar granularity. Panel B avoids forward-filled return observations but contains variable one- to six-day intervals. It is not a fixed-horizon return panel and does not solve intraday synchronization. Panel A maintains nominal business-day spacing but introduces carry-forward zeros. The two-panel design exposes this trade-off rather than eliminating it.

The fourth is stress-window definition. The study uses prespecified COVID-19 and 2024 windows and three finite thresholds. Different stress definitions or thresholds could produce different point estimates. Such alternatives should be treated as new exploratory analyses rather than retrofitted after seeing the current results.

The fifth is scope. Earlier monthly transmission regressions were withdrawn after the macroeconomic data audit identified invalid CPI observations in the extracted series. The present paper therefore does not make claims about commodity-to-inflation, export or monthly exchange-rate transmission.

\subsection{Discussion summary}

The main contribution of the results is not a single copula estimate or a single stress p-value. It is the separation of conclusions by how well they survive competing sources of specification uncertainty.

The absence of multiplicity-adjusted stress discoveries survives both calendar constructions and all three stress definitions. Complete rejection of the eight tested static copulas survives both calendars for Gold-Brent, Gold-WTI and Brent-WTI. The equivalent conclusion does not survive for Cocoa-Gold, Cocoa-Brent or Cocoa-WTI, whose adequacy results are calendar-sensitive. Commodity-Cedi non-rejection is also stable across calendars, but it must be interpreted alongside the unresolved Cedi marginal.

This pattern supports a broader methodological lesson: robustness in dependence analysis cannot be inferred from one fitted model, one preprocessing rule or one nominally significant cell. In this setting, the strongest conclusions are the ones that remain after calendar construction, marginal adequacy, absolute goodness of fit and multiple testing are all made explicit.

\section{Robustness and Limitations}
\label{sec:robustness_and_limitations}

\subsection{Calendar robustness as the main specification check}

The principal robustness exercise is built into the study design rather than added after the main analysis. Panel A preserves the documented business-day construction with limited forward filling, while Panel B uses only consecutive common observation dates. The two panels are not treated as a primary analysis and a secondary robustness check. They are co-equal representations of the observation process.

This comparison yields three different outcomes. The absence of Bonferroni- and Benjamini-Hochberg-adjusted stress discoveries is stable across both panels for the combined, COVID-only and 2024-only stress families. Complete rejection of all eight tested static copulas is also stable for Gold-Brent, Gold-WTI and Brent-WTI. By contrast, the copula adequacy conclusions for Cocoa-Gold, Cocoa-Brent and Cocoa-WTI are not stable: all eight families are rejected in Panel A, while only a subset is rejected in Panel B.

The value of the second panel is therefore not that it confirms every result. It reveals where a substantive conclusion depends on the calendar construction. This is why the manuscript distinguishes calendar-robust findings from calendar-sensitive findings rather than reporting a single preferred preprocessing choice.

\subsection{Marginal convergence and admissibility audit}

The marginal stage was re-audited after the original analysis to ensure that optimizer convergence was treated as a hard admissibility condition. All 20 Panel A candidate fits were checked under the rule that a model is valid only when the optimizer converges and the log likelihood, estimated parameters, standardized residuals, conditional volatility and PIT values are finite.

One Panel A candidate, Cedi AR(1)-EGARCH-t, failed the convergence gate. Excluding it from both the adequacy pool and the AIC fallback did not change any of the previously selected Panel A models. Cocoa and Gold remain AR(1)-GJR-GARCH-t, Brent remains AR(1)-GJR-GARCH-skew-t, WTI remains AR(1)-EGARCH-t, and the Cedi remains an inadequate AR(1)-GJR-GARCH-t reference.

The Panel B marginal rebuild applies the same rule. Cocoa, Gold, Brent and WTI have an adequate selected candidate, while the Cedi again has no adequate tested specification and is represented only by an inadequate reference model. The persistence of the Cedi problem across both calendars is therefore not explained by a failed optimizer in one branch of the analysis.

\subsection{Copula-estimator implementation audit}

Because Panel A and Panel B were produced through different reconstruction branches, the copula GOF implementation was checked directly before cross-panel conclusions were frozen. The audit compared the estimators used by the historical Panel A GOF path with those used by the Panel B publication rebuild.

For all eight families and all ten pair cells, the compared fitting routes produced the same observed fitted parameters and the same observed GOF statistics when supplied with the same PIT inputs. The audit also confirmed that the bootstrap refits the copula in every replicate and that the current publication p-value convention is

\[
\hat p=\frac{\text{exceedances}+1}{B+1}.
\]

No material estimator mismatch was detected. A new Panel A GOF bootstrap was therefore not required on implementation-equivalence grounds. The observed cross-panel differences in GOF classification are attributed to the different panel inputs and marginal selections rather than to different copula estimators.

This audit does not prove that the GOF procedure is the only appropriate test. It establishes only that the cross-panel comparison is not confounded by an unnoticed change in the implemented estimator.

\subsection{PIT clipping and ties}

The stored PIT values are clipped to \([10^{-6},1-10^{-6}]\) before the copula rank transformation. Panel B contains no exact PIT ties. Panel A contains one three-way Cedi tie at the lower clipping boundary, on 19 September 2022, 21 February 2023 and 29 December 2023.

Using the saved standardized residuals, the underlying order of the three observations was recovered. Their correct lower-tail ranks are 3, 2 and 1, respectively, rather than the common average rank of 2 induced by clipping. Only two pseudo-observations change, and the maximum displacement is approximately \(3.32\times10^{-4}\).

The correction does not change membership at \(q=0.025\), \(0.05\) or \(0.10\). It therefore does not change which observations enter the prespecified finite-tail sets. Across the 32 Panel A commodity-Cedi GOF cells, the largest observed change in the GOF statistic is \(0.000262\). For the borderline Cocoa-Cedi Frank cell, the observed statistic is unchanged to the reported precision.

The PIT tie is therefore a documented numerical artifact, but not a material driver of the current stress or GOF conclusions. A full bootstrap rerun was not warranted on this basis.

\subsection{Alternative Cedi treatments}

Panel A contains an additional seven-treatment Cedi sensitivity analysis designed to examine whether nominal commodity-Cedi stress patterns depend on how the zero-heavy Cedi process is handled. The treatments include the baseline GARCH reference, raw-rank treatment, AR-only filtering, weekly aggregation, a hurdle Student-t treatment, a hurdle-mixture treatment and a hurdle-Markov treatment.

These analyses are diagnostic rather than co-equal headline specifications. Each treatment evaluates 12 commodity-Cedi diagnostic cells without multiplicity correction. The number of nominal 5\% findings varies across treatments:

\begin{table}[htbp]
\centering
\small
\begin{tabularx}{\textwidth}{lX}
\toprule
\textbf{Cedi treatment} & \textbf{Nominal findings out of 12} \\
\midrule
Baseline GARCH reference & 3 \\
No-GARCH raw ranks & 3 \\
AR-only & 2 \\
Weekly frequency & 0 \\
Hurdle Student-t & 2 \\
Hurdle mixture & 2 \\
Hurdle Markov & 2 \\
\bottomrule
\end{tabularx}
\end{table}

A negative Brent-Cedi lower-tail difference at \(q=0.05\) appears repeatedly under several daily treatments but disappears under the weekly treatment. This lack of invariance is one reason the pattern is not interpreted as a structural commodity-Cedi result.

The synchronized Panel B seven-treatment matrix was not run automatically. The follow-up rule was fixed in advance: it would be triggered only if at least one Panel B commodity-Cedi stress cell survived Bonferroni or Benjamini-Hochberg correction. No such cell survived under any of the three stress definitions, so the trigger was not activated. The Panel A seven-treatment exercise therefore remains supporting diagnostic evidence rather than a two-panel headline result.

\subsection{Extreme-value sensitivity}

The main POT/GPD threshold is fixed at the 90th percentile of the standardized-residual loss series. For the Panel B Cedi reference, the fitted shape parameter remains positive over the examined threshold range from 0.85 to 0.975, approximately 0.63 to 0.79.

This threshold stability is informative about the fitted reference series, but it does not resolve the marginal-model problem. Because the Cedi's selected marginal specification fails the PIT adequacy test, the EVT quantities remain conditional on that reference residual process. The threshold sensitivity exercise therefore supports robustness of the fitted heavy-tail pattern within that model, not identification of a model-free structural tail index.

\subsection{Withdrawn monthly transmission branch}

An earlier version of the project included monthly regressions linking aggregated tail measures to Cedi depreciation, inflation and export growth. That branch is excluded from the present manuscript.

The reason is data validity, not an unfavorable statistical result. The audit of the extracted Bank of Ghana macroeconomic series identified invalid literal zero values in headline inflation for part of 2023. The processed export series also contained a source gap that had previously been masked by the extraction pipeline. Because the monthly regressions depend on those inputs, their coefficients and p-values are withdrawn from the current evidence base.

The paper therefore makes no empirical claim about monthly commodity-to-Cedi, inflation or export transmission. Reintroducing that question would require reconstruction of the macroeconomic panel from independently verified authoritative data followed by a fresh analysis.

This withdrawal is an important part of the reproducibility record. A historical result is not retained simply because it had already appeared in an earlier draft.

\subsection{Remaining statistical limitations}

Several limitations remain after the robustness work.

First, the Cedi marginal remains unresolved. The tested continuous conditional distributions do not adequately reproduce its zero-heavy daily process. This limitation propagates into Cedi-related copula and EVT quantities.

Second, the copula analysis is bivariate and static. Rejection of all eight tested families for some commodity pairs suggests the need for richer dependence structures, but the present analysis does not determine whether a dynamic, regime-switching, mixture, vine or nonparametric alternative would be adequate.

Third, neither calendar construction achieves perfect market synchronization. Panel A can create carry-forward zero returns, while Panel B uses intervals ranging from one to six calendar days. Intraday timing differences are not modeled.

Fourth, the stress analysis is conditional on the prespecified COVID-19 and 2024 windows, the thresholds \(q=0.025,0.05,0.10\), a moving-block length of 20 and the chosen multiplicity procedures. The study does not claim that every reasonable alternative definition would produce the same point estimates or p-values.

Fifth, static GOF non-rejection for the commodity-Cedi pairs may reflect limited ability to discriminate among weak dependence structures. The present study does not perform a separate empirical power analysis for those real-data cells, so low power is not asserted as a demonstrated explanation.

Finally, the commodity price files are described only to the level established by the audited sources and processing scripts. The manuscript does not make unverified claims about continuous-futures roll rules or back adjustment.

\subsection{Robustness summary}

The robustness work changes the paper in a substantive way. It confirms some findings, weakens others and removes one entire analysis branch.

The stress conclusion survives the two calendar constructions and the prespecified multiplicity procedures. Complete eight-family copula rejection survives for Gold-Brent, Gold-WTI and Brent-WTI but not for the three cocoa-related commodity pairs. The marginal convergence audit leaves the Panel A selections unchanged but confirms that the Cedi remains unresolved. The GOF implementation audit rules out an estimator mismatch as the source of the cross-panel differences. The PIT clipping audit identifies a real numerical artifact but shows that it is too small to alter the current conclusions. The Cedi sensitivity analysis shows that nominal daily patterns are not invariant to alternative treatment and frequency. The monthly transmission branch is removed because its input data failed audit.

The resulting manuscript is therefore narrower than the historical draft, but the remaining claims are supported by a clearer chain from source data to statistical conclusion.

\section{Conclusion}
\label{sec:conclusion}

This study examined commodity tail dependence and exchange-rate risk in Ghana using two co-equal calendar constructions, explicit marginal-model adequacy checks, eight static bivariate copula families, finite-threshold tail-concentration measures and multiplicity-controlled stress inference. The purpose was not to identify a single preferred preprocessing rule or a single best-fitting copula, but to determine which conclusions remain defensible after calendar alignment, marginal specification, absolute model adequacy and multiple testing are treated explicitly.

Three findings survive that scrutiny. First, the stress-period conclusion is calendar-robust. Across the combined COVID-19 plus 2024 definition, COVID-19 alone and 2024 alone, neither the business-day panel nor the common-observation panel produces a Bonferroni- or Benjamini-Hochberg-adjusted discovery among the 60 finite-tail comparisons in each family. This result does not establish that dependence was unchanged during those episodes. It means that the tested stress-versus-calm differences do not provide multiplicity-adjusted evidence strong enough to support individual changes under the prespecified thresholds and bootstrap design.

Second, absolute copula adequacy is itself sensitive to the way the observation calendar is constructed. Gold-Brent, Gold-WTI and Brent-WTI reject all eight tested static families in both panels, giving a calendar-robust indication that the candidate set is inadequate for these pairs. The same conclusion does not hold for Cocoa-Gold, Cocoa-Brent and Cocoa-WTI. Each rejects all eight families under the business-day construction but only a subset under common-observation alignment. These cocoa-related results show that calendar choice can alter higher-order dependence conclusions even when the corresponding return series remain strongly correlated across panels. For applied dependence analysis, calendar construction is therefore part of the statistical specification rather than a neutral preprocessing detail.

Third, the Cedi remains the main unresolved marginal problem. None of the tested continuous marginal candidates passes the full adequacy gate under either calendar construction. Commodity-Cedi copula non-rejections are stable across panels, but they should not be interpreted as proof that the tested families are correct. Likewise, the Cedi EVT results indicate a heavy fitted loss tail under the retained reference model, but those quantities remain conditional on marginal misspecification. The separate clipping audit shows that the identified Panel A PIT tie is numerically minor and does not drive the substantive conclusions. The larger issue is the inability of the tested continuous marginal models to reproduce the zero-heavy daily Cedi process adequately.

The paper's contribution is therefore methodological as much as empirical. It shows how apparently simple conclusions about commodity dependence can change once absolute fit, calendar construction and multiple testing are considered jointly. It also demonstrates the value of retaining negative and inconvenient findings rather than forcing a single narrative. Some conclusions survive both panels, some become calendar-sensitive, one historical analysis branch is withdrawn after its data fail audit, and Cedi-related evidence remains explicitly conditional. This produces a narrower set of claims than the earlier manuscript, but a more defensible one.

For Ghana, the results should not be read as evidence that commodity markets are unimportant for the Cedi or that crisis-period dependence cannot increase. The analysis addresses a specific daily-source-frequency tail question and does not estimate causal pass-through, fiscal effects, reserve implications or lower-frequency macroeconomic transmission. Those questions remain open. Future work would benefit from a separately validated macroeconomic data pipeline, more flexible models for the Cedi's mixed discrete-continuous behavior, and dynamic, regime-switching, vine or nonparametric dependence models for commodity pairs that fail the static adequacy tests.

The broader lesson is straightforward. In multivariate risk analysis, robustness is not demonstrated by one statistically significant cell, one preferred calendar or one lowest-AIC model. The most credible conclusions are those that survive alternative defensible data constructions, explicit model-adequacy checks and the inferential burden created by multiple comparisons.


\section*{Data and Code Availability}
The analysis code and publication-rebuild materials are maintained in the project repository:
\url{https://github.com/Jonathanerrils/tail-dependence-ghana}. The reconstructed manuscript is restricted to claims supported by the audited data, code and reproducible outputs documented in the project.

\bibliographystyle{plainnat}
\bibliography{references}

\end{document}
